\section{Frontend}
\label{sec:frontend}
\noindent Dans un premier temps, nous avons développé un site Web, présenté dans la Section~\ref{sec:web}.
Une fois que la première version du site était prête, nous avons commencé le développement
d'une application mobile, présentée dans la Section~\ref{sec:mobile},
facilitant l'utilisation sur smartphone.

\subsection{Site Web}
\label{sec:web}
\noindent L'application Web a été développée en utilisant React \cite{React} et Vite \cite{Vite}.
Vite permet d'automatiser la création de l'architecture du projet, le serveur live de développement et la compilation en build production.\\

\noindent Le site Web est composé d'une page d'accueil présentant le projet et invitant les utilisateurs à installer l'application mobile.
Il y a aussi une page de connexion, nécessitant une adresse mail et un mot de passe.
Si l'utilisateur n'a pas encore de compte, il peut en créer un via la page de création de compte.
Il y a, de plus, une page consacrée à la recherche d'itinéraires (Section~\ref{sec:webrecherche}) et une page de profil utilisateur (Section~\ref{sec:webprofil}).

\subsubsection{Page de recherche d'itinéraires}
\label{sec:webrecherche}
\noindent La page principale de l'application est celle permettant la recherche d'itinéraire (Figure~\ref{fig:webrecherche}).

\begin{figure}[H]
    \centering
    \includegraphics[width=\textwidth]{screenshots/site/PageItineraire.png}
    \caption{Capture d'écran de la page de recherche d'itinéraire.}
    \label{fig:webrecherche}
\end{figure}

\noindent La partie gauche de cette page contient le menu de recherche, permettant d'entrer les adresses de départ et d'arrivée du trajet, l'heure maximale d'arrivée ou la durée maximale de trajet, ainsi que le type de vélo utilisé.
Les adresses sont autocomplétées grâce à l'utilisation de l'API Base Adresse Nationale (BAN \cite{BAN}), proposée par le gouvernement français.
En appuyant sur le bouton "Calculer les itinéraires", l'utilisateur envoie la requête de recherche d'itinéraires à l'API.
Les trajets proposés sont ensuite affichés sur la carte, et leur description apparaît dans la partie gauche de l'écran.\\

\noindent Un trajet est décrit par son type : rapide, sécurisé ou un compromis entre les deux.
Nous affichons la dénivelée positive et négative rencontrée lors du trajet, la distance et le temps de parcours estimé.
Celui-ci dépend du profil sportif de l'utilisateur.
Nous affichons aussi le pourcentage de pistes cyclables composant le trajet, le pourcentage de zone où la vitesse des véhicules est limitée à 30 km/h ainsi que le pourcentage de voies éclairées.\\

\noindent De plus, l'utilisateur a accès à un petit graphique représentant l'évolution de la dénivelée pendant le trajet.
En cliquant sur le bouton d'information situé à côté du type de trajet, un texte récapitulant les caractéristiques de sécurité du trajet apparaît.\\

\noindent De plus, la page de recherche d'itinéraire propose une carte.
Sur celle-ci s'affichent les différents trajets suggérés.
Un cours descriptif du trajet apparaît quand la souris le survole.
Un clique sur le trajet provoque la sélection du trajet et l'apparition du descriptif détaillé du trajet dans la partie gauche de l'écran.\\

\noindent Sur la carte se trouvent trois boutons principaux.

\noindent Le bouton "Trafic en temps réel" déclenche l'apparition des points de trafic ainsi que de leur légende.
Les données du trafic sont récupérées grâce à l'API Avatar \cite{Avatar} proposée par le Cerema.
Ces données sont utilisées avec une heure et demie de décalage et sont actualisées toutes les cinq minutes.\\

\noindent On trouve aussi sur la carte le bouton "Ajouter un signalement".
Une fois ce bouton cliqué, l'utilisateur peut cliquer à un endroit sur la carte.
Il peut alors créer un signalement, visible par les autres utilisateurs et pris en compte dans le calcul des itinéraires.
L'utilisateur peut sélectionner le type de signalement parmi "Accident", "Travaux", "Danger" et "Obstacle", et entrer s'il le souhaite une description de l'évènement.\\

\noindent Le bouton "Calques" permet de modifier le fond de carte.
L'utilisateur peut donc l'ajuster selon ses goûts.
Il peut choisir entre une carte sombre, une carte avec le relief, les rues, la vue satellite...

\subsubsection{Page profil utilisateur}
\label{sec:webprofil}
\noindent La seconde page du site Web est la page présentant le profil de l'utilisateur (Figure~\ref{fig:webprofil}).

\begin{figure}[H]
    \centering
    \includegraphics[width=\textwidth]{screenshots/site/PageUtilisateur.png}
    \caption{Capture d'écran de la page de profil de l'utilisateur.}
    \label{fig:webprofil}
\end{figure}

\noindent À partir de cette page, l'utilisateur peut accéder à ses informations.
Un clic sur le bouton "Modifier mon compte" provoque l'ouverture d'une fenêtre modale permettant à l'utilisateur de modifier son nom, son adresse mail, sa date de naissance, son mot de passe et son niveau sportif (Débutant, Intermédiaire ou Avancé).

\noindent L'utilisateur a aussi accès à ses adresses de préférence, qui lui sont recommandées lors de la recherche d'itinéraire pour permettre une utilisation quotidienne rapide.
Il peut modifier ses adresses via le bouton "Modifier mes adresses".

\noindent L'utilisateur peut aussi enregistrer ses vélos, attribuer à chacun un nom personnalisé, un type (Ville, VTT ou Route) et indiquer s'il est ou non électrique.
Les vélos de l'utilisateur lui sont proposés lors de la recherche d'itinéraire.\\

\noindent Sur la page de profil se trouve l'historique de l'utilisateur, ainsi que ses statistiques, directement calculées par rapport à l'historique.
L'utilisateur a aussi le choix de modifier son historique.
Lors d'un clic sur un trajet de celui-ci, une fenêtre modale s'ouvre, lui permettant d'accéder à toutes les informations du trajet effectué (Figure~\ref{fig:historic}).

\begin{figure}[H]
    \centering
    \includegraphics[width=\textwidth]{screenshots/site/HistoricModal.png}
    \caption{Capture d'écran de la fenêtre modale d'historique.}
    \label{fig:historic}
\end{figure}

\subsubsection{Référencement et SEO}
\noindent Afin que le site puisse apparaître dans les résultats de recherche, nous avons rendu accessible un fichier \texttt{robots.txt} autorisant l'indexation et un fichier \texttt{sitemap.xml} qui présente les différentes pages du site devant être indexées avec un poids d'importance et leur url.
Nous avons également enregistré le site auprès de la Google Search Console \cite{GoogleSearchConsole} et de Bing Webmaster Tools \cite{BingWebmasterTools} pour demander une indexation explicitement et suivre les statistiques d'apparition du site dans les résultats de recherche.
Combiner ces deux plateformes permet l'apparition dans la plupart des moteurs de recherches et également dans les résultats des modèles IA.

\noindent Pour améliorer le score de ranking nous avons optimisé les temps de chargement et mis des balises \texttt{<meta>} dynamiques en fonction des pages. Nous avons également essayé de suivre au maximum les bonnes pratiques de structuration de pages.

\noindent Actuellement le référencement n'est qu'à son début et la page n'est pas encore bien indexée. En effet, le nom de domaine vient d'être déposé et le site vient d'être mis en ligne, donc le score de pertinence de la page est encore bas. Pour faire augmenter la crédibilité du site auprès des moteurs de recherche, il faudrait que des pages sérieuses et reconnues parlent de ce projet et contiennent un lien renvoyant sur le site.

\subsection{Application mobile}
\label{sec:mobile}

\noindent L'application mobile a été développée en utilisant React Native \cite{ReactNative}, permettant de créer une application compatible à la fois avec Android et iOS. Nous avons utilisé le framework Expo \cite{Expo} pour faciliter le développement et le déploiement de l'application.\\

\noindent L'application mobile propose les mêmes fonctionnalités que le site web, mais de manière réarrangée pour mieux correspondre à l'utilisation sur smartphone. L'application embarque également la navigation en temps réel avec instructions détaillée dans la Sous-section~\ref{sec:mobile-navigation}.

\noindent Pour le moment l'application n'est pas publiée sur les stores d'applications.

\noindent L'application est composée principalement d'une page de recherche d'itinéraires (page principale) et d'une page de profil utilisateur qui sont accessibles via un menu en bas de l'écran. Elle est également composée de pages annexes de formulaires comme la page de connexion, de création de compte, de modification de profil, d'ajout de vélo ou encore de détail de l'historique. Les pages principales de l'application sont présentées dans la Figure~\ref{fig:app_screens}.

\begin{figure}[H]
    \centering
    \begin{subfigure}[b]{0.2\textwidth}
        \centering
        \includegraphics[width=\textwidth]{screenshots/application/apercu_itineraire}
        \caption{Aperçu des itinéraires}
        \label{fig:apercu_itineraire}
    \end{subfigure}
    \hfill
    \begin{subfigure}[b]{0.2\textwidth}
        \centering
        \includegraphics[width=\textwidth]{screenshots/application/detail_itineraire}
        \caption{Détail d'un itinéraire}
        \label{fig:detail_itineraire}
    \end{subfigure}
    \hfill
    \begin{subfigure}[b]{0.2\textwidth}
        \centering
        \includegraphics[width=\textwidth]{screenshots/application/navigation}
        \caption{Navigation en temps réel}
        \label{fig:navigation}
    \end{subfigure}
    \hfill
    \begin{subfigure}[b]{0.2\textwidth}
        \centering
        \includegraphics[width=\textwidth]{screenshots/application/profile}
        \caption{Page de profil utilisateur}
        \label{fig:profil}
    \end{subfigure}
    \caption{Différentes vues de l'application mobile.}
    \label{fig:app_screens}
\end{figure}

\subsubsection{Page de recherche d'itinéraires}
\label{sec:mobile-recherche}

\noindent La page de recherche d'itinéraires est la page principale de l'application mobile. Elle est présentée dans la Sous-Figure~\ref{fig:apercu_itineraire}. Au démarrage de l'application, l'utilisateur est directement sur cette page pour pouvoir faire rapidement une recherche d'itinéraire.\\

\noindent En haut de la page se trouve le bloc de recherche qui contient les champs d'adresses de départ et d'arrivée. Un bouton flèche permet d'agrandir le bloc de recherche pour faire apparaître la sélection du vélo et de la contrainte horraire. Les adresses sont autocomplétées comme sur le site Web. La spécificité de l'application sur la recherche est la possibilité de partir de la position GPS actuelle. Dans la plupart des cas d'utilisations classiques sur mobile, l'utilisateur part de sa position pour se rendre quelque part, il est donc important que ce choix soit facilement accessible.\\

\noindent La carte occupe la totalité de l'écran en arrière plan. Elle démarre sur la position actuelle. L'utilisateur peut interagir avec la carte, la déplacer et zoomer. Un bouton permet de recentrer la vue sur la position. La vue se centre automatiquement sur le point qui vient d'être entré comme adresse de départ ou d'arrivée, ou sur les deux si les deux adresses sont complétées.\\

\noindent Lorsque les trajets sont calculés, ils s'affichent sur la carte. C'est cet état qui est montré dans la Sous-Figure~\ref{fig:apercu_itineraire}. L'utilisateur peut cliquer sur un trajet pour le sélectionner. Un panneau sur le bas de l'écran apparaît et permet également de sélectionner un itinéraire. Ces panneaux affichent le type d'itinéraire et quelques informations supplémentaires. Pour voir tout le détail de l'itinéraire, un bouton permet d'ouvrir une fenêtre modale, montrée dans la Sous-Figure~\ref{fig:detail_itineraire}. L'utilisation d'une modale permet que les panneaux ne soient pas trop grand, et l'utilisateur peut donc visualiser la carte en même temps que les détails de son itinéraire.\\

\noindent Un bouton permet de changer le fond de carte, comme sur le site. Le thème de la carte peut être clair ou sombre selon le choix de l'utilisateur et est par défaut calqué sur le thème du téléphone.\\

\subsubsection{Navigation en temps réel}
\label{sec:mobile-navigation}

\noindent Le mode navigation est le point de différence principal de l'application mobile et ce qui en fait son réel intérêt par rapport au site. L'utilisateur peut entrer dans le mode navigation une fois un itinéraire sélectionné. Une fois entré dans le mode, l'interface s'épure en enlevant le panneau de recherche, la carte s'incline légèrement pour mieux montrer la route et un nouveau panneau apparaît. Cette interface est présentée dans la Sous-Figure~\ref{fig:navigation}.
Ce nouveau panneau affiche les instructions de navigation en temps réel. Ces instructions sont mises à jour en temps réel en fonction de la position de l'utilisateur et sont calculées côté backend.
L'avantage de faire de cette façon est que le backend a la connaissance complète de la route (intersection, rond-points, etc), le frontend lui ne voit qu'une ligne de trajet. Une icone indique le type d'action que doit faire l'utilisateur à la prochaine intersection (tourner, continuer tout droit, prendre une sortie de rond-point, etc).
Pour faciliter l'expérience du cycliste, l'écran du smartphone reste allumé tant qu'il est en mode navigation. Nous avons également implémenté le guidage vocal, qui utilise la librairie Text To Speech du téléphone.\\
